\section{Related Work}

\subsection{Causal Inference in Recommendation Systems}

Causal inference models have been increasingly explored in the context of recommendation systems, primarily due to their ability to uncover and leverage cause-effect relationships within user interaction data. Past research [\cite{Pearl2009Causality}, \cite{Shmueli2010Explanatory}, \cite{Spirtes2000Causal}, \cite{Peters2017Elements}] has highlighted the significance of employing Structural Equation Models (SEMs) and causal diagrams to identify and model these relationships. Such methodologies aid in enhancing user profiling by isolating pivotal factors influencing user decisions. Additionally, recent studies [\cite{Louizos2017Causal}, \cite{Yao2021Debiased}] have explored the integration of causal analysis with machine learning techniques to improve the interpretability and fairness of recommendations.

Our work builds on these principles by integrating advanced causal inference methods with adaptive learning. Unlike traditional models, our approach employs SEMs and causal diagrams not only to extract causal features but also to guide adaptive learning processes, ensuring refined exploration-exploitation strategies that enhance recommendation precision.

\subsection{Adaptive Learning and Exploration-Exploitation Trade-off}

The exploration-exploitation trade-off in adaptive learning has been a focus of research in recommender systems. Multi-armed bandit (MAB) frameworks, such as those discussed in [\cite{Gittins2011Multi}, \cite{Auer2002Using}, \cite{Vermorel2005Multi}, \cite{Russo2018Tutorial}, \cite{Bubeck2012Regret}], are widely utilized to optimize recommendation strategies through balancing exploration (trying new recommendations) and exploitation (leveraging known preferences). Variants like Thompson Sampling and epsilon-greedy strategies have been particularly effective in real-time learning environments. Studies [\cite{Zhou2015Survey}, \cite{Li2010Contextual}] further demonstrate how contextual information can enhance adaptive learning algorithms by tailoring recommendations more closely to user-specific contexts.

In our proposed method, we extend these adaptive learning strategies by incorporating causal insights from user data. This combination allows us to inform action decisions more accurately, improving personalization and recommendation relevance while preserving a stable exploration-exploitation balance.

\subsection{Semantic Content Mapping and NLP Techniques}

Semantic content mapping has received considerable attention, particularly in leveraging NLP methodologies for enhancing the contextual relevance of recommendations. Techniques such as Latent Dirichlet Allocation (LDA) and entity recognition have been utilized extensively in prior works [\cite{Blei2003LDA}, \cite{Manning2014Stanford}, \cite{Pennington2014GloVe}, \cite{Devlin2019BERT}]. These strategies facilitate aligning recommendations with user preferences by deciphering thematic structures and contextual nuances in user data. Further, advancements in contextual embeddings and dynamic topic models [\cite{Liu2019Text}, \cite{Pang2008Opinion}] contribute to refining semantic alignment between content and users' interests.

Our approach advances these methodologies by integrating them with adaptive learning algorithms. By using advanced exploration-exploitation strategies alongside semantic mapping, we achieve a high degree of personalization, ensuring that recommendations are not only thematic but also dynamically resonate with evolving user contexts and preferences.

```

```bibtex
